\section{Conclusion}
The goal was to find out whether the Reachy Mini can detect a person,
recognise their activity and give expressive feedback within the limits
of its own CPU. For the technical part of that goal the answer is
The complete pipeline runs on the robot's Raspberry Pi CM4
without a GPU or a laptop. Pose estimation takes about 34 ms per
frame, features and forest a few milliseconds, and the camera's 10~FPS,
not the computation, sets the pace. The same code and models run in
simulation, so they are developed on a laptop and deployed unchanged.
Recognised actions are answered in real time with antenna signatures and
head and body following.

The answer to ``how well'' is more limited. In live use, wave, clapping
and jumping jacks were reliable enough to demonstrate, and squats were
recognised with about one repetition in seven missed. Push-ups were not.
Scores on public datasets overstate what happens in a room: the
six-class model recognises waves with $F_1 = 0.918$ on dataset subjects
but $F_1 = 0.581$ on our robot recordings. For the CMAS items, pick-up
held up in a room ($F_1 = 0.933$). Chair rise and arms above the head
were recognised live on the robot, but they have not yet been measured
on a correctly performed room recording. The robot's position also
matters: it must see the whole person, and a waist-up view costs
recognition.

The prototype, its models, the training and export tools, and the coach
were realised. A CMAS score, the stool-step item and clinical validation
were not; they were outside the scope. The ``engaging'' part of the
research question was not evaluated, since no user study was done. The
prototype is therefore a research platform: it shows that
monitoring at home with this robot is technically feasible, but it
cannot yet be used to monitor children with JDM.

That said, there are four recommended next steps:
\begin{enumerate}
    \item Record labelled sessions in real rooms through the robot's own
    camera, following the CMAS instructions, and use them as the reference
    for room performance.
    \item Collect data for the classes that lack it: push-ups, clapping
    on robot footage, and the stool step.
    \item Run a small user study with children and families to measure
    engagement, which the research question assumes, but this project did
    not test.
    \item If CMAS scoring becomes the goal, define the scoring criteria
    with the WKZ specialists and validate the system against their scores.
\end{enumerate}